\chapter{Hogyan gondolkodjunk egy Velran alkalmazásról?}

\noindent\textbf{Tanulási út: Gyors út: alap.} Építsd fel azt a mentális modellt, amelyre minden későbbi fejezet épül.\par\medskip

\rustbridgeblockhu{A structok, enumok, \code{Option}, \code{Result} és explicit függvényparaméterek továbbra is ismerős typed építőelemek.}{A model compiler számára láthatóvá teszi az alkalmazás adatstruktúráját, a capability pedig explicitté teszi a külső erőforrásokhoz tartozó authorityt.}{Nem kell új objektummodellt megtanulnod. Maradj a typed értékeknél és explicit bemeneteknél; azt tanuld meg, hol tesz köréjük a Velran security-határt.}


Route vagy adatbázis előtt ezt kérdezd: \emph{milyen üzleti adat létezik, milyen alakban, milyen műveletek végezhetők rajta, és milyen határok között?} A Velran több biztonsági és helyességi garanciája abból származik, hogy ezek az adatalakok, inputok, query-k és erőforrás-hozzáférések explicit szerződések.

\section{Először a fogalmak, utána az oldalak}

Egy terméklistánál indulj a \emph{Product} fogalomból, annak typed alakjából, az olvasó/módosító műveletekből, a szükséges capability-kből és policy-kből, és csak ezután a HTTP felületből. A gondolkodás iránya nagyjából:

\begin{lstlisting}
uzleti fogalom
    -> typed adatalak
        -> muveletek
            -> capability-k es policy-k
                -> HTTP route
                    -> HTML / JSON
\end{lstlisting}

Ez a sorrend segít, hogy az üzleti szabályok ne ad hoc route-okba, template-ekbe vagy összefűzött SQL-be szóródjanak.

\section{Mit jelent a \code{model}?}

A könyvben két közeli fogalommal találkozunk:

\begin{itemize}
  \item \textbf{adatmodell vagy domain modell}: annak emberi leírása, milyen üzleti fogalmak és kapcsolatok léteznek az alkalmazásban;
  \item \textbf{\code{model} deklaráció}: a Velran konkrét, fordító által ismert adatalakja.
\end{itemize}

Például:

\begin{lstlisting}[language=Velran]
model Product {
    id: i64
    name: String
    price: i64
}
\end{lstlisting}

Ez azt mondja a fordítónak, hogy egy \code{Product} rekordnak pontosan ilyen mezői és típusai vannak. A typed SQL query eredményalakja ehhez a szerződéshez igazodik, és ugyanez az érték később HTML- vagy JSON-kimenetben is használható.

\subsection*{A \code{model} nem automatikus adatbázisséma}

A fenti deklaráció önmagában nem hoz létre \code{products} táblát. Az adatbázissémát migration hozza létre, például:

\begin{lstlisting}[language=SQL]
CREATE TABLE products (
    id INTEGER PRIMARY KEY,
    name TEXT NOT NULL,
    price INTEGER NOT NULL
);
\end{lstlisting}

A rétegek kapcsolódnak, de külön felelősségük van: a migration kezeli a tartós sémát, a \code{model} az alkalmazás typed adatszerződése, a query pedig explicit módon képez sorokat modellé. A model ezért nem hoz létre automatikusan adatbázis-entitást vagy azonos nevű táblát.

\subsection*{A \code{model} nem mutable objektum}

A \code{Product} nem klasszikus objektumorientált osztály. Nincs mögötte rejtett \code{save()} metódus, implicit adatbázis-kapcsolat vagy mutable objektuméletciklus. Ha adatot akarunk módosítani, explicit műveletet és explicit erőforrást használunk.

Itt még az aláírás a lényeg: a \code{listProducts} \code{Db} olvasási hozzáférést kap, a \code{createProduct} pedig explicit \code{Transaction} erőforrást a módosításhoz. A teljes SQL-törzseket a 3. fejezet mutatja meg, amikor már a teljes request útvonalat látjuk.

\section{A model a rétegek közti szerződés}

A \code{model} a rétegek között mozgó typed szerződés.

\begin{lstlisting}
Database row
    |  typed query
    v
Product model
    |-------------------|
    v                   v
HTML page            JSON API
    |
    v
felhasznalo
\end{lstlisting}

Ha egy query \code{Product}-ot ígér, a visszaadott mezőknek ehhez az alakhoz kell igazodniuk. Ugyanez a typed rekord-szerződés tér vissza később az object authorizationnél és a JSON-kimenetnél.

\section{Az üzleti fogalomhoz műveletek tartoznak}

Az üzleti adathoz műveletek tartoznak. Ilyen a \code{Product.list} vagy az \code{Article.publish}. Kis programban ezek lehetnek top-level függvények; nagyobb programban a később bemutatott \code{object} namespace csoportosíthatja őket. Fontos különbség:

\begin{center}
\begin{tabularx}{\textwidth}{@{}lY@{}}
\toprule
Elem & Kérdés, amelyre válaszol \\
\midrule
\code{model} & Milyen alakú adat a \code{Product}? \\
\callable{query} & Hogyan olvassuk vagy módosítjuk a tartós adatot? \\
\code{object} & Mely üzleti fogalomhoz tartoznak ezek a műveletek? \\
\callable{page}/\callable{action} & Hogyan kapcsolódik az üzleti művelet a webes handlerhez? \\
\code{route} & Milyen HTTP útvonalon, inputtal és policy-val érhető el? \\
\bottomrule
\end{tabularx}
\end{center}

\section{Capability: a hozzáférés is része a modellnek}

Velranban nem csak az adat típusa explicit, hanem az is, milyen erőforráshoz fér hozzá egy művelet. Egy függvény nem ``valahonnan'' kap adatbázist, fájlrendszert vagy current usert.

\begin{lstlisting}[language=Velran]
#[page]
fn home(ctx: PageContext, db: Db)
    -> Result<Html, PageError> {
    let products = listProducts(db)?;
    return Ok(html {
        <ul>
            @for product in products {
                <li>{{ product.name }}</li>
            }
        </ul>
    });
}
\end{lstlisting}

A \code{db: Db} az aláírásban látható. Fájlkezelésnél AppFs capability, módosító adatbázis-műveletnél Transaction, autentikált műveletnél pedig az auth policy és a hozzá tartozó typed kontextus válik relevánssá.

Tervezési szabály: ha egy műveletnek érzékeny erőforrás kell, annak a kódból és a policy-ból látszania kell.

\section{A route nem az üzleti modell}

A route fontos, de a rendszer külső határa. Például:

\begin{lstlisting}[language=Velran]
route create POST "/products"
    form name<String> price<i64>
    validate name length 1 100 price range 0 100000000
    public => create;
\end{lstlisting}

A route a külső HTTP-határt írja le: metódus, path, typed input, validáció és policy. Maga a \code{Product} ettől függetlenül HTML-, JSON-, admin- vagy más route mögött is újrahasználható.

Ezért nagyobb alkalmazásban különösen hasznos a következő határ:

\begin{lstlisting}
domain/model     : mi az adat es milyen muvelet ertelmes rajta?
HTTP route       : ki, hogyan es milyen policy-val hivhatja?
runtime config   : milyen infrastruktura-es eroforras-hatarok kozott futhat?
\end{lstlisting}

\section{Egy egyszerű tervezési módszer}

Amikor új feature-t kezdesz, a következő öt kérdés jó kiindulópont:

\begin{enumerate}
  \item \textbf{Mi az üzleti főnév?} -- Product, Article, Invoice, WikiPage.
  \item \textbf{Milyen typed adatra van ténylegesen szükség?} -- ne automatikusan a teljes adatbázissort másold át.
  \item \textbf{Mely műveletek olvasnak, és melyek módosítanak?}
  \item \textbf{Milyen capability kell hozzájuk?} -- Db, Transaction, AppFs, auth és named hálózati egress integration.
  \item \textbf{Milyen HTTP/policy felületen tesszük elérhetővé?} -- route, input, validáció, auth, cache, rate limit.
\end{enumerate}

CRUD-nál ez pár perc; nagyobb rendszerben megakadályozza, hogy a route-ok, SQL-ek és jogosultsági szabályok szétcsússzanak.

\section{Mit vigyél tovább a következő fejezetekbe?}

Egy szabályt vigyél tovább: a typed adat, query, tranzakció, capability és HTTP/security policy maradjon explicit, fölösleges objektumhierarchia nélkül. A következő fejezetek telepítik a rendszert, majd ezt a modellt egy teljes \code{Product} requesten alkalmazzák.

\section{Gyakorlat: rajzold fel a határtérképet}
\paragraph{Gyakorlási mód: vezetett.} Használd az előszóban meghatározott vezetett módszert.

Válassz egy egyszerű, jól ismert funkciót, például jegyzetet, feladatot vagy készletelemet. Kódírás előtt sorold fel a megbízható állapotot, a kérésből érkező adatokat, a tartós rekordokat, a módosítási pontokat és a kívülről látható válaszokat. Ezután jelöld meg, mely átmenetekhez kell validáció, autorizáció, tranzakció vagy explicit capability.

A cél a rejtett feltételezések láthatóvá tétele még a kód előtt.

\section{Tipikus hiba: a route-okból indulunk ki a modell helyett}
A route-központú tervezés gyakran olyan handlereket eredményez, amelyek újra és újra ugyanazokat az üzleti szabályokat fedezik fel. Ha több route ugyanazt az állapotot parsolja, ugyanazt a tulajdonosi ellenőrzést végzi vagy ugyanazokat az invariánsokat építi újra, vidd vissza a gondolkodást a típusos modellekhez, query-khez és explicit módosítási határokhoz.

\section{Folyamatos mintaprojekt: Secure Notes}
A compilerrel ellenőrzött \path{examples/secure-notes/checkpoints/01-baseline/main.vrn} \emph{Secure Notes} legyen a folyamatos mintaprojekt. A baseline authenticated, principal szerint szűrt hozzáférést, bounded inputot, CSRF-védett formokat, ownership ellenőrzést és tranzakciós módosításokat kényszerít ki; a későbbi checkpointok ezeket az invariánsokat nem gyengítik.
